\paragraph{window-\texorpdfstring{$6$}{6} 的三原子 moment 测度、Hankel 刚性与 involutive thermodynamics}
在 window-\(6\) 的完整有限尺度谱、bin-fold 纤维群胚代数的 Wedderburn 分解、输出熵账本以及三能级 escort 模型已经固定之后，还可把稳定类型、微态、碰撞、矩阵代数与有限热力学进一步压缩为单一的三原子 Stieltjes 对象。记
\[
d_6(w):=d_6^{\mathrm{bin}}(w),
\qquad
\mu_6:=\sum_{w\in X_6}\delta_{d_6(w)}=8\delta_2+4\delta_3+9\delta_4,
\]
并沿用本文关于 involutive groupoid algebra 的记号
\[
\mathcal A_6:=\CC[\mathcal R_6^{\mathrm{bin}}],
\qquad
\mathcal A_6^{\alpha,\pm}:=\{a\in \mathcal A_6:\ \alpha(a)=\pm a\}.
\]
下列结果表明，window-\(6\) 的多层对象并非平行排布：\(\mu_6\) 的矩序列具有三阶湮没多项式，其 Hankel 秩恰等于支撑点数，其前三个 moment 在 \((\mathcal A_6,\alpha)\) 中分别实现为中心、反不动扇区与全代数，而输出 Rényi 谱、dyadic 可逆容量与 gauge 型热力学量又全部由同一测度通过显式积分恢复。

\begin{theorem}[window-\texorpdfstring{$6$}{6} 纤维多重度测度的湮没三次式与有理母函数]\label{thm:xi-time-part9z-window6-threeatom-moment-ogf}
\leanverified{paper\_xi\_time\_part9z\_window6\_threeatom\_moment\_ogf}
定义矩序列
\[
M_n:=\int t^n\,\dd\mu_6(t)
=
\sum_{w\in X_6}d_6(w)^n
=
8\cdot 2^n+4\cdot 3^n+9\cdot 4^n
\qquad (n\ge 0).
\]
则其 ordinary generating function 满足
\[
\mathcal M_6(z):=\sum_{n\ge 0}M_n z^n
=
\frac{21-125z+182z^2}{(1-2z)(1-3z)(1-4z)}.
\]
特别地，\((M_n)_{n\ge 0}\) 满足精确线性递推
\[
M_{n+3}=9M_{n+2}-26M_{n+1}+24M_n
\qquad (n\ge 0),
\]
其对应特征多项式为
\[
p_6(\lambda):=(\lambda-2)(\lambda-3)(\lambda-4)
=
\lambda^3-9\lambda^2+26\lambda-24.
\]
\end{theorem}
\begin{proof}
由 \(\mu_6=8\delta_2+4\delta_3+9\delta_4\)，对每个 \(n\ge 0\) 直接有
\[
M_n=8\cdot 2^n+4\cdot 3^n+9\cdot 4^n.
\]
因此
\[
\mathcal M_6(z)
=
\frac{8}{1-2z}+\frac{4}{1-3z}+\frac{9}{1-4z},
\]
通分后得到
\[
\mathcal M_6(z)
=
\frac{21-125z+182z^2}{(1-2z)(1-3z)(1-4z)}.
\]
又
\[
(1-2z)(1-3z)(1-4z)=1-9z+26z^2-24z^3,
\]
故生成函数分母对应的递推正是
\[
M_{n+3}-9M_{n+2}+26M_{n+1}-24M_n=0.
\]
把 \(\lambda\) 视为移位算子的特征值变量，即得所示特征多项式。证毕。
\end{proof}

\begin{theorem}[window-\texorpdfstring{$6$}{6} 三原子 moment 测度的 Hankel 平坦化]\label{thm:xi-time-part9z-window6-threeatom-hankel-flatness}
对 \(r\ge 0\)，记
\[
H_r^{(6)}:=(M_{i+j})_{0\le i,j\le r}.
\]
则有
\[
\det H_2^{(6)}=1152>0,
\qquad
\det H_3^{(6)}=0.
\]
因此 \((M_n)\) 的 Hankel 秩恰为 \(3\)。等价地，\(\mu_6\) 恰为三原子测度，而其支撑由定理 \ref{thm:xi-time-part9z-window6-threeatom-moment-ogf} 的湮没多项式唯一恢复为
\[
\operatorname{supp}\mu_6=\{2,3,4\}.
\]
\leanverified{paper\_xi\_time\_part9z\_window6\_threeatom\_hankel\_flatness}
\end{theorem}
\begin{proof}
对每个 \(r\ge 0\)，记
\[
V_r:=
\begin{pmatrix}
1&1&1\\
2&3&4\\
\vdots&\vdots&\vdots\\
2^r&3^r&4^r
\end{pmatrix},
\qquad
D:=\operatorname{diag}(8,4,9).
\]
由
\[
M_{i+j}=8\cdot 2^{i+j}+4\cdot 3^{i+j}+9\cdot 4^{i+j},
\]
可得 Hankel 分解
\[
H_r^{(6)}=V_r D V_r^{\mathsf T}.
\]
于是对所有 \(r\) 都有 \(\operatorname{rank}H_r^{(6)}\le 3\)，从而
\[
\det H_3^{(6)}=0.
\]

当 \(r=2\) 时，\(V_2\) 为 \(3\times 3\) Vandermonde 矩阵，故
\[
\det H_2^{(6)}
=
(\det V_2)^2\det D
=
\bigl((3-2)(4-2)(4-3)\bigr)^2\cdot 8\cdot 4\cdot 9
=
1152.
\]
因此 \(\operatorname{rank}H_2^{(6)}=3\)，遂 \((M_n)\) 的 Hankel 秩恰为 \(3\)。

另一方面，三原子测度的矩序列最小递推阶恰等于其 Hankel 秩；故定理 \ref{thm:xi-time-part9z-window6-threeatom-moment-ogf} 中的三次湮没多项式已达到最小可能阶数。其根 \(2,3,4\) 两两不同，因此正是 \(\mu_6\) 的支撑点。证毕。
\end{proof}

\begin{theorem}[window-\texorpdfstring{$6$}{6} 的 \texorpdfstring{$0$}{0}--\texorpdfstring{$1$}{1}--\texorpdfstring{$2$}{2} moment 在 involutive algebra 中的实现]\label{thm:xi-time-part9z-window6-moment-involutive-algebra-realization}
window-\(6\) 的前三个纤维 moment 满足
\[
M_0=\dim Z(\mathcal A_6)=21,
\]
\[
M_1=\dim \mathcal A_6^{\alpha,-}=64,
\]
\[
M_2=\dim \mathcal A_6=212.
\]
因此，稳定类型数、微态总数与有序碰撞对总数并非三套彼此平行的量，而分别是同一测度 \(\mu_6\) 的 \(0\)、\(1\)、\(2\) 阶 moment，并在 involutive groupoid algebra 中依次实现为中心、反不动扇区与全代数。
\end{theorem}
\begin{proof}
由命题 \ref{prop:fold-bin-groupoid-algebra-wedderburn} 与推论 \ref{cor:fold-bin-m6-fiber-hist}，
\[
\mathcal A_6
\cong
M_2(\CC)^{\oplus 8}\oplus M_3(\CC)^{\oplus 4}\oplus M_4(\CC)^{\oplus 9}.
\]
于是
\[
\dim Z(\mathcal A_6)=8+4+9=21=M_0,
\]
且
\[
\dim \mathcal A_6
=
8\cdot 2^2+4\cdot 3^2+9\cdot 4^2
=
212
=
M_2.
\]

对 \(M_1\)，一方面，纤维分割覆盖全部 \(2^6\) 个微态，故
\[
M_1=\sum_{w\in X_6}d_6(w)=2^6=64.
\]
另一方面，本文已审计数据给出
\[
\dim \mathcal A_6^{\alpha,-}=64.
\]
遂得第二式。证毕。
\end{proof}

\begin{corollary}[window-\texorpdfstring{$6$}{6} 的 \texorpdfstring{$\alpha$}{alpha}-固定扇区恰计数非对角碰撞]\label{cor:xi-time-part9z-window6-fixed-sector-distinct-collision}
有精确恒等式
\[
\dim \mathcal A_6^{\alpha,+}
=
M_2-M_1
=
\sum_{w\in X_6}d_6(w)\bigl(d_6(w)-1\bigr)
=
148.
\]
换言之，\(\alpha\)-固定扇区并不计数全部有序碰撞对，而恰计数有序互异碰撞对；全部对角配对 \((\omega,\omega)\) 的总数 \(64\) 则完全落在 \(\alpha\)-反不动扇区。
\end{corollary}
\begin{proof}
由于 \(\alpha\) 为对合，\(\mathcal A_6=\mathcal A_6^{\alpha,+}\oplus \mathcal A_6^{\alpha,-}\)。故由定理 \ref{thm:xi-time-part9z-window6-moment-involutive-algebra-realization}，
\[
\dim \mathcal A_6^{\alpha,+}
=
\dim \mathcal A_6-\dim \mathcal A_6^{\alpha,-}
=
212-64
=
148.
\]
另一方面，
\[
\sum_{w\in X_6}d_6(w)\bigl(d_6(w)-1\bigr)
=
\sum_{w\in X_6}d_6(w)^2-\sum_{w\in X_6}d_6(w)
=
M_2-M_1.
\]
合并两式即得结论。证毕。
\end{proof}

\begin{theorem}[window-\texorpdfstring{$6$}{6} 的前三个 moment 唯一恢复完整纤维直方图]\label{thm:xi-time-part9z-window6-histogram-from-three-moments}
记
\[
N_2:=\#\{w\in X_6:\ d_6(w)=2\},
\qquad
N_3:=\#\{w\in X_6:\ d_6(w)=3\},
\qquad
N_4:=\#\{w\in X_6:\ d_6(w)=4\}.
\]
则 \((N_2,N_3,N_4)\) 可由 \((M_0,M_1,M_2)\) 唯一恢复，且显式公式为
\[
N_2=6M_0-\frac72 M_1+\frac12 M_2,
\]
\[
N_3=-8M_0+6M_1-M_2,
\]
\[
N_4=3M_0-\frac52 M_1+\frac12 M_2.
\]
代入
\[
(M_0,M_1,M_2)=(21,64,212)
\]
即得
\[
(N_2,N_3,N_4)=(8,4,9).
\]
\leanverified{paper\_xi\_time\_part9z\_window6\_histogram\_from\_three\_moments}
\end{theorem}
\begin{proof}
按定义，
\[
N_2+N_3+N_4=M_0,
\]
\[
2N_2+3N_3+4N_4=M_1,
\]
\[
4N_2+9N_3+16N_4=M_2.
\]
由前两式得
\[
N_3+2N_4=M_1-2M_0.
\]
再用第三式减去第二式的两倍，得到
\[
N_3+4N_4=M_2-2M_1.
\]
两式相减即得
\[
2N_4=M_2-5M_1+6M_0,
\]
故
\[
N_4=3M_0-\frac52M_1+\frac12M_2.
\]
代回
\(
N_3+2N_4=M_1-2M_0
\)
得
\[
N_3=-8M_0+6M_1-M_2.
\]
最后由 \(N_2=M_0-N_3-N_4\)，得到
\[
N_2=6M_0-\frac72M_1+\frac12M_2.
\]
代入 \((M_0,M_1,M_2)=(21,64,212)\) 便得 \((8,4,9)\)。证毕。
\end{proof}

\begin{theorem}[window-\texorpdfstring{$6$}{6} 输出分布的精确 Rényi 谱由三原子测度闭合]\label{thm:xi-time-part9z-window6-threeatom-renyi-spectrum}
令
\[
p_6(w):=\frac{d_6(w)}{64}
\qquad (w\in X_6)
\]
为 bin-fold 在稳定类型层上的输出分布。则对任意 \(q\in\RR\setminus\{1\}\)，其 Rényi 熵满足
\[
H_q(p_6)
=
\frac{1}{1-q}
\Bigl[
\log\bigl(8\cdot 2^q+4\cdot 3^q+9\cdot 4^q\bigr)-6q\log 2
\Bigr].
\]
相应极限值为
\[
H_0(p_6)=\log 21,
\]
\[
H_1(p_6)=\frac{37}{8}\log 2-\frac{3}{16}\log 3,
\]
\[
H_2(p_6)=10\log 2-\log 53,
\]
\[
H_\infty(p_6)=4\log 2.
\]
\end{theorem}
\begin{proof}
由 \(\mu_6=8\delta_2+4\delta_3+9\delta_4\)，有
\[
\sum_{w\in X_6}p_6(w)^q
=
\frac{8\cdot 2^q+4\cdot 3^q+9\cdot 4^q}{64^q}.
\]
代入 Rényi 熵定义即得第一式。

\(q\to 0\) 时，有效支撑大小为 \(|X_6|=21\)，故
\[
H_0(p_6)=\log 21.
\]
\(q\to 1\) 时，Shannon 熵为
\[
H_1(p_6)
=
-\sum_{w\in X_6}p_6(w)\log p_6(w)
=
6\log 2-\frac{1}{64}\sum_{w\in X_6}d_6(w)\log d_6(w),
\]
从而
\[
H_1(p_6)
=
6\log 2-\frac{1}{64}\Bigl(8\cdot 2\log 2+4\cdot 3\log 3+9\cdot 4\log 4\Bigr)
=
\frac{37}{8}\log 2-\frac{3}{16}\log 3.
\]

对 \(q=2\)，有
\[
\sum_{w\in X_6}p_6(w)^2
=
\frac{M_2}{64^2}
=
\frac{212}{4096}
=
\frac{53}{1024},
\]
故
\[
H_2(p_6)=-\log\frac{53}{1024}=10\log 2-\log 53.
\]
最后，
\[
\max_{w\in X_6}p_6(w)=\frac{4}{64}=\frac{1}{16},
\]
故
\[
H_\infty(p_6)=-\log\frac{1}{16}=4\log 2.
\]
证毕。
\end{proof}

\begin{theorem}[window-\texorpdfstring{$6$}{6} escort 壳层族的射影可识别性与四室相图]\label{thm:xi-time-part9z-window6-escort-projective-identifiability}
\leanverified{paper\_xi\_time\_part9z\_window6\_escort\_projective\_identifiability}
对任意 \(q\in\RR\)，记
\[
Z_6(q):=8\cdot 2^q+4\cdot 3^q+9\cdot 4^q,
\]
并定义三条壳层 escort 概率
\[
\pi_2(q):=\frac{8\cdot 2^q}{Z_6(q)},
\qquad
\pi_3(q):=\frac{4\cdot 3^q}{Z_6(q)},
\qquad
\pi_4(q):=\frac{9\cdot 4^q}{Z_6(q)}.
\]
则参数 \(q\) 可由单一交叉比精确恢复：
\[
q=
\frac{
\log\!\Bigl(\dfrac{9\,\pi_3(q)^2}{2\,\pi_2(q)\pi_4(q)}\Bigr)
}{\log(9/8)}.
\]
因此，window-\(6\) 的 escort 壳层族不是仅凭数值拟合辨识的一维曲线，而是严格射影可识别的参数族。

此外，三壳层权重的次序随 \(q\) 恰发生三次交换；相应临界温度为
\[
q_{3=4}:=\frac{\log(4/9)}{\log(4/3)}\approx -2.8188416793,
\]
\[
q_{2=4}:=\frac{\log(8/9)}{\log 2}\approx -0.1699250014,
\]
\[
q_{2=3}:=\frac{\log 2}{\log(3/2)}\approx 1.7095112914.
\]
从而三壳层相图严格分裂为四个 chamber：
\[
q<q_{3=4}:\qquad \pi_2(q)>\pi_3(q)>\pi_4(q),
\]
\[
q_{3=4}<q<q_{2=4}:\qquad \pi_2(q)>\pi_4(q)>\pi_3(q),
\]
\[
q_{2=4}<q<q_{2=3}:\qquad \pi_4(q)>\pi_2(q)>\pi_3(q),
\]
\[
q>q_{2=3}:\qquad \pi_4(q)>\pi_3(q)>\pi_2(q).
\]
\end{theorem}
\begin{proof}
直接计算得
\[
\frac{\pi_3(q)^2}{\pi_2(q)\pi_4(q)}
=
\frac{(4\cdot 3^q)^2}{(8\cdot 2^q)(9\cdot 4^q)}
=
\frac{2}{9}\left(\frac{9}{8}\right)^q,
\]
整理即得
\[
q=
\frac{
\log\!\bigl(9\pi_3(q)^2/(2\pi_2(q)\pi_4(q))\bigr)
}{\log(9/8)}.
\]
故 \(q\) 已由单一交叉比唯一恢复。

再看壳层次序。三条成对比值满足
\[
\frac{\pi_2(q)}{\pi_3(q)}=2\left(\frac{2}{3}\right)^q,
\qquad
\frac{\pi_2(q)}{\pi_4(q)}=\frac{8}{9}\,2^{-q},
\qquad
\frac{\pi_3(q)}{\pi_4(q)}=\frac{4}{9}\left(\frac{3}{4}\right)^q.
\]
它们都是 \(q\) 的严格递减函数，因此每一对壳层只会发生一次次序交换。令各比值等于 \(1\) 即得三条临界曲线
\[
8\cdot 2^q=4\cdot 3^q,
\qquad
8\cdot 2^q=9\cdot 4^q,
\qquad
4\cdot 3^q=9\cdot 4^q,
\]
解得 \(q_{2=3}\)、\(q_{2=4}\)、\(q_{3=4}\) 三式。

由其数值顺序
\[
q_{3=4}<q_{2=4}<q_{2=3}
\]
可知：当 \(q<q_{3=4}\) 时，三条比值均大于 \(1\)，故 \(\pi_2>\pi_3>\pi_4\)；越过 \(q_{3=4}\) 后，只有 \(\pi_3/\pi_4\) 变为小于 \(1\)，故 \(\pi_2>\pi_4>\pi_3\)；越过 \(q_{2=4}\) 后，再有 \(\pi_2/\pi_4<1\)，故 \(\pi_4>\pi_2>\pi_3\)；最后越过 \(q_{2=3}\) 后，\(\pi_2/\pi_3<1\)，故 \(\pi_4>\pi_3>\pi_2\)。证毕。
\end{proof}

\begin{theorem}[window-\texorpdfstring{$6$}{6} 可逆容量的 Stieltjes 势表示与 Mellin--Stieltjes 对偶]\label{thm:xi-time-part9z-window6-stieltjes-capacity-duality}
对整数 \(B\ge 0\)，定义 dyadic 预言机下的最大可逆恢复容量
\[
C_6(B):=\sum_{w\in X_6}\min(2^B,d_6(w)).
\]
则有精确表示
\[
C_6(B)=\int \min(2^B,t)\,\dd\mu_6(t).
\]
若再定义尾函数
\[
T_6(s):=\mu_6([s,\infty)),
\]
则
\[
C_6(B)=\int_0^{2^B}T_6(s)\,\dd s.
\]
进一步，把整数矩沿实参数延拓为
\[
M_s:=\int t^s\,\dd\mu_6(t)
=
8\cdot 2^s+4\cdot 3^s+9\cdot 4^s
\qquad (s>0),
\]
则有精确 Mellin--Stieltjes 恒等式
\[
M_s=s\int_0^\infty u^{s-1}T_6(u)\,\dd u
\qquad (s>0).
\]
对 window-\(6\) 而言，
\[
T_6(s)=
\begin{cases}
21,& 0<s\le 2,\\
13,& 2<s\le 3,\\
9,& 3<s\le 4,\\
0,& s>4,
\end{cases}
\]
从而
\[
C_6(0)=21,\qquad
C_6(1)=42,\qquad
C_6(B)=64\ \ (B\ge 2).
\]
\end{theorem}
\begin{proof}
由 \(\mu_6=\sum_{w\in X_6}\delta_{d_6(w)}\)，直接有
\[
\int \min(2^B,t)\,\dd\mu_6(t)
=
\sum_{w\in X_6}\min(2^B,d_6(w))
=
C_6(B).
\]

对任意 \(a>0\)，恒等式
\[
\min(a,t)=\int_0^a \mathbf 1_{\{t\ge s\}}\,\dd s
\]
成立。将其对 \(\mu_6\) 积分，得到
\[
\int \min(a,t)\,\dd\mu_6(t)
=
\int_0^a \mu_6([s,\infty))\,\dd s
=
\int_0^a T_6(s)\,\dd s.
\]
取 \(a=2^B\) 即得容量公式。

对 \(s>0\)，对有限正测度施行 Stieltjes 分部积分，得到
\[
\int t^s\,\dd\mu_6(t)
=
s\int_0^\infty u^{s-1}\mu_6([u,\infty))\,\dd u
=
s\int_0^\infty u^{s-1}T_6(u)\,\dd u.
\]
最后，由 \(\mu_6=8\delta_2+4\delta_3+9\delta_4\) 可立即读出 \(T_6\) 的分段表达；逐段积分便得
\[
C_6(0)=\int_0^1 21\,\dd s=21,
\qquad
C_6(1)=\int_0^2 21\,\dd s=42,
\]
而当 \(B\ge 2\) 时，\(2^B\ge 4\)，故
\[
C_6(B)
=
\int_0^4 T_6(s)\,\dd s
=
21\cdot 2+13\cdot 1+9\cdot 1
=
64.
\]
证毕。
\end{proof}
\endinput
